\documentclass[11pt]{article}
\usepackage[T1]{fontenc}
\usepackage[utf8]{inputenc}
\usepackage{lmodern}
\usepackage[margin=1in]{geometry}
\usepackage{amsmath,amssymb,amsthm,mathtools}
\usepackage{microtype}
\usepackage{enumitem}
\usepackage{booktabs}
\usepackage{xcolor}
\usepackage[colorlinks=true,linkcolor=blue!55!black,citecolor=blue!55!black,
  urlcolor=blue!55!black,bookmarksnumbered=true]{hyperref}
\hypersetup{pdftitle={Sinkless Orientation in Incremental dynamic-LOCAL},
  pdfsubject={Self-contained proofs of deterministic and randomized locality upper bounds}}
\newtheorem{theorem}{Theorem}[section]
\newtheorem{lemma}[theorem]{Lemma}
\newtheorem{corollary}[theorem]{Corollary}
\newtheorem{proposition}[theorem]{Proposition}
\theoremstyle{definition}
\newtheorem{definition}[theorem]{Definition}
\newtheorem{remark}[theorem]{Remark}
\newcommand{\dist}{\operatorname{dist}}
\newcommand{\diam}{\operatorname{diam}}
\newcommand{\outdeg}{d^+}
\newcommand{\SO}{\textsc{SO}}
\newcommand{\Dyn}{\textnormal{dynamic-LOCAL}}
\newcommand{\On}{\textnormal{online-LOCAL}}
\newcommand{\Ball}{B}
\newcommand{\Prb}{\mathbb P}
\newcommand{\EE}{\mathbb E}
\setlist[enumerate]{itemsep=3pt,topsep=5pt}
\setlist[itemize]{itemsep=3pt,topsep=5pt}
\setlength{\emergencystretch}{2em}
\setcounter{tocdepth}{1}

\title{Sinkless Orientation in Incremental \Dyn\\[3pt]
  \large Double-logarithmic deterministic and triple-logarithmic randomized locality}
\author{}
\date{25 October 2026}

\begin{document}
\maketitle

\begin{abstract}
We maintain an orientation of an incrementally growing graph in which every
vertex of degree at least three has an outgoing edge. We prove worst-case
locality $O(\log\log n)$ deterministically and
$O(\log\log\log n)$ with randomization against an oblivious insertion
adversary. In the randomized result, for any fixed $c>0$, one event of
probability at least $1-n^{-c}$ guarantees correctness after \emph{every}
update; the radius bound holds on every random tape.

The main argument is a problem-specific transfer from online to dynamic
locality. First reduce monotonically to maximum degree three. Query a single
online algorithm only at vertices whose surrounding neighborhoods have
become permanent, and freeze their incident directions. Every remaining
vertex is close to a vertex of degree at most two. Orienting a short path
toward such a vertex repairs a sink without disturbing frozen outputs or
creating another sink. An online radius $r$ gives dynamic radius $2r+4$,
with no loss in success probability for an oblivious sequence.

For completeness, we also prove the static ingredients: the deterministic
clustering-and-greedy algorithm and the randomized marking-and-shattering
algorithm underlying the known sequential bounds. No algorithmic theorem
from the cited literature is needed as a black box in the upper-bound proofs.
\end{abstract}

\tableofcontents
\clearpage

\section{Problem, models, and statements}\label{sec:model}

All input graphs are finite, simple, undirected, and have at most a known
number $n\ge2$ of vertices. Vertices have distinct identifiers, which are
used only for deterministic choices. Graph distance is unweighted; the
distance between distinct components is infinite. Write $B_G(v,t)$ for the
closed radius-$t$ ball. An orientation assigns one direction to each edge.
A \emph{defect} is a vertex of degree at least three and outdegree zero.
A \emph{sinkless orientation} is an orientation without defects. Thus
vertices of degree zero, one, or two are exempt. This convention makes the
problem well-defined throughout a graph's growth, including on forests.

\subsection{Incremental dynamic locality}

Initially the graph is empty. An update either inserts one edge between
existing vertices or inserts one vertex together with any set of edges from
it to existing vertices. Edges and vertices are never deleted, and parallel
edges are not inserted. The algorithm must output an orientation after every
update. Intermediate computations during an update are not outputs.

Global computation and persistent internal memory are unrestricted. Locality
restricts changes to the \emph{output}, not access to the input or running
time. We use the endpoint encoding: a vertex reports the direction of every
incident edge, and the reports at its two endpoints must agree. Reversing an
old edge therefore changes both endpoints' outputs.

For precision, after insertion of $uv$ our radius-$R$ algorithms change old
edge directions only if both endpoints of the changed edge lie in
\begin{equation}\label{eq:affected}
 B_G(u,R)\cap B_G(v,R),
\end{equation}
where $G$ is the graph \emph{after} the insertion. The directions of newly
inserted edges are also assigned in this region. For insertion of a vertex
$z$, the corresponding allowed region is $B_G(z,R)$. These are the usual
affected regions in the induced-neighborhood formulation of incremental
\Dyn\ \cite{models}. In proofs it is convenient to first bound distance to
the \emph{nearer} endpoint. The elementary inclusion
\begin{equation}\label{eq:union-intersection}
 B_G(u,a)\cup B_G(v,a)
 \ \subseteq\ B_G(u,a+1)\cap B_G(v,a+1)
\end{equation}
converts that bound to~\eqref{eq:affected}.

\begin{definition}[Randomized guarantee]\label{def:random}
A randomized incremental algorithm has radius $R(n)$ and error
$\varepsilon(n)$ against an \emph{oblivious adversary} if:
\begin{enumerate}
\item its output changes obey the radius bound for every random tape; and
\item for every entire insertion sequence fixed independently of its private
random bits, the probability that its outputs are legal after all updates
is at least $1-\varepsilon(n)$.
\end{enumerate}
The sequence includes vertex identifiers and any ordering data supplied with
an update. The algorithm receives updates one at a time and does not know
the future sequence.
\end{definition}

This is a Monte Carlo guarantee with a hard locality bound. It is stronger
than a separate success guarantee for an individual update. Our randomized
theorem does not assume that future updates can depend on observed random
outputs. The deterministic theorem permits such adaptive choices.

\begin{theorem}[Main results]\label{thm:main}
Sinkless orientation on general graphs admits the following incremental
\Dyn\ algorithms:
\begin{enumerate}[label=(\roman*)]
\item a deterministic algorithm of worst-case locality $O(\log\log n)$;
\item for every fixed $c>0$, a randomized algorithm of worst-case locality
$O(\log\log\log n)$ and error at most $n^{-c}$ in
Definition~\ref{def:random}.
\end{enumerate}
Both statements allow edge insertions and vertex insertions with incident
edges. Constants in the randomized locality may depend on $c$.
\end{theorem}

Iterated logarithms in asymptotic bounds are truncated below by one, so the
statements include small values of $n$. Unless written $\ln$, logarithms are
base two. These bounds concern radius per update, not the number of changed
edges, total computation, or amortized work.

\subsection{The online intermediate model}

In \On, the input graph $H$ is fixed but initially unknown. Vertices are
queried one at a time in an arbitrary order. On query $v$, an algorithm of
radius $r$ sees the subgraph induced by the union of all radius-$r$ balls
of vertices queried so far, with their identifiers and any port labels. It
has unrestricted persistent global memory. It must give $v$'s final output,
which can never be revised. For orientation, that output specifies every
incident edge direction. At the end, all reports must agree and give a
sinkless orientation. All algorithms below use $r\ge1$; an extra layer can
always be included when degree or boundary-incidence information is needed.

A randomized online algorithm has error $\varepsilon$ if, for every fixed
graph and fixed query order independent of its coins, its entire final output
is legal with probability at least $1-\varepsilon$. Its radius bound holds
on every tape. These definitions follow the locality models
of~\cite{models,quantum}. In particular, information or provisional plans
kept in private memory need not coincide with the final outputs already
requested from vertices.

\subsection{Proof plan and provenance}\label{sec:provenance}

The dynamic argument is given first, in
Sections~\ref{sec:reduction}--\ref{sec:random-transfer}. It transforms an
online radius $r$ into dynamic radius $2r+4$. Then
Sections~\ref{sec:static-tools}--\ref{sec:static-random} construct the required
online algorithms and thereby complete Theorem~\ref{thm:main}.

The static bounds are established results. Ghaffari, Harris, and Kuhn
proved the deterministic and randomized sequential bounds
\cite[Theorems 7.2 and 7.3]{GHK}. Our deterministic construction follows the
particularly simple greedy and clustering proof of Balliu et
al.~\cite[Section 3]{simple}. The randomized construction follows the
marking-and-shattering method of Ghaffari and Su~\cite{GS} and its sequential
adaptation in~\cite[Algorithm 3]{GHK}. The composition of sequential phases
goes back to Ghaffari, Kuhn, and Maus~\cite{GKM}; we give the easier
global-memory online simulation explicitly. We include proofs of all of
these ingredients, with the graph classes and error conventions needed here.

The additional ingredient is to delay irrevocable decisions until input
neighborhoods have stabilized. This reverses, for this particular problem,
the usual direction of simulation from dynamic to online locality
\cite{models}. Short local repair is related to the viewpoint of local
mending~\cite{mending}; here the algorithm maintains a particular family of
repairable states rather than extending arbitrary partial solutions.
Related repairs toward low-degree vertices or cycles appear in Nolin's
recent simple distributed sinkless-orientation algorithm~\cite{Nolin},
whose complexity bounds concern distributed communication rounds.

Recent work of Bourreau et al.~\cite{coloring} obtains optimal logarithmic
dynamic locality for bipartite 3-coloring and discusses the difficulty of
using online decisions when the graph can still change. Sinkless orientation
has a useful special feature: in a subcubic graph, a degree-three vertex is
saturated forever, while a nonsaturated vertex is exempt from the output
constraint. Our proof uses both properties.

There is a relevant historical qualification. Earlier versions of
\cite{quantum} claimed a general dynamic derandomization theorem that would
also imply the deterministic numerical bound here. That claim was withdrawn
in the corrected version~4, Section~3.7 and the errata (13 February 2026).
The proof below is independent of that claim. We make no general
online-to-dynamic or dynamic derandomization assertion.

\section{A monotone reduction to maximum degree three}\label{sec:reduction}

\begin{lemma}\label{lem:degree-reduction}
An incremental sinkless-orientation algorithm on subcubic graphs yields one
on general graphs with the same vertex bound, success probability, and
locality, apart from a minimum radius of one. The reduction is deterministic
and depends only on the insertion sequence.
\end{lemma}
\begin{proof}
For each vertex, permanently select its first three arriving incident edges,
or all of them if fewer have arrived. In a simultaneous arrival, use a fixed
identifier order. Let $H$ contain exactly those edges selected by both
endpoints. The graph $H$ has maximum degree three and only grows: the status
of an edge is decided when it arrives and never changes.

Run the subcubic algorithm on $H$. Orient an edge selected by exactly one
endpoint out of that endpoint, permanently. Orient an edge selected by
neither endpoint according to any fixed rule. Use the algorithm's directions
for the edges of $H$.

If $\deg_G(v)\ge3$, then $v$ has selected three edges. If one of them was
selected only by $v$, it is an outgoing edge. Otherwise all three belong to
$H$, so $\deg_H(v)=3$ and the orientation of $H$ gives $v$ an outgoing edge.
Thus correctness in $H$ implies correctness in $G$ at every time.

An original edge insertion either inserts that same edge into $H$ or leaves
$H$ unchanged. Since $H\subseteq G$, distances in $G$ are no larger than
distances in $H$, so every permitted change in $H$ is permitted in $G$.
All other old edges keep their directions. There are no new vertices. The
entire support-graph sequence is fixed when the original insertion sequence
is fixed, so the probability guarantee is preserved as well.
\end{proof}

We prove the subcubic algorithm for single-edge insertions on a fixed set of
$n$ initially isolated virtual vertices. Actual vertex arrivals can be mapped
to unused virtual identifiers in $[n]$ in arrival order; physical identifiers
need not be used internally. If ports are needed, assign them permanently
in support-edge arrival order. Unused vertices remain isolated and impose
no constraint. Section~\ref{sec:vertices} explains why
inserting a vertex with several edges preserves the same locality bound.

\section{The deterministic online-to-dynamic transfer}\label{sec:transfer}

\begin{theorem}\label{thm:transfer}
Suppose $A$ is a deterministic \On\ algorithm for sinkless orientation on
subcubic graphs with $n$ vertices and radius $r=r(n)\ge1$. There is a
deterministic incremental \Dyn\ algorithm on the same graph class with
radius at most $2r+4$.
\end{theorem}

\subsection{Free and committed vertices}

Set $\tau=r+1$. In the current subcubic graph $H$, define
\begin{equation}\label{eq:sets}
 F=\{v:\deg_H(v)\le2\},\qquad
 \delta(v)=\dist_H(v,F),\qquad
 C=\{v:\delta(v)>\tau\},\qquad U=V(H)\setminus C.
\end{equation}
If a component has no vertex in $F$, its distances to $F$ are infinite.
Vertices in $F$ are \emph{free}, and vertices in $C$ are \emph{committed}.
We run one persistent instance of $A$, querying a vertex exactly when it
first joins $C$. Its incident edge directions are then frozen. Vertices in
$U$ have temporary directions, which are never passed to $A$ as if they
were outputs of $A$.

\begin{lemma}[Monotonicity]\label{lem:monotonicity}
Under subcubic edge insertions, $\delta(v)$ never decreases. Hence $C$ only
grows, and each vertex is queried at most once.
\end{lemma}
\begin{proof}
Immediately before inserting $ab$, both $a$ and $b$ have degree at most two,
so both are in the old free set. The new free set is a subset of the old
one; only endpoints whose degrees reach three can leave it.

Consider a shortest new path from $x$ to a new free vertex. If it does not
use $ab$, it was already an old path to an old free vertex. If it uses $ab$,
its prefix up to the first endpoint of $ab$ is an old path to an old free
vertex. In either case the old distance to the free set is at most the
length of the new path. If the new distance is infinite there is nothing
to prove.
\end{proof}

\begin{lemma}[Permanence of revealed input]\label{lem:permanence}
When $x$ enters $C$, its radius-$r$ online view, including one additional
boundary layer, is permanent. At every later time the complete transcript
of queries to $A$ is a valid online prefix on the current graph, and also
on every future graph in the insertion sequence.
\end{lemma}
\begin{proof}
At commitment, every vertex within distance $r+1$ of $x$ has degree three.
No edge can subsequently be inserted incident to any of those vertices.
The old radius-$r$ ball can therefore neither lose an edge nor acquire a
new incident edge.

Every later path from $x$ of length at most $r+1$ uses only edges present at
commitment time. Otherwise its first new edge has an old-edge prefix of
length at most $r$, so its first endpoint was already saturated. That edge
could not have been inserted. Consequently no new shortcut changes either
ball or its boundary data.

Finally, all vertices in all previously revealed radius-$r$ balls are
saturated. Consequently no subsequent insertion changes an edge induced
by their union or their incidences. Edges revealed when a further ball is
added are consistent with the same current or future graph. The identifiers,
ports, graph views, and query order are therefore exactly those of an online
prefix on that fixed graph. This proves the assertion for the union-of-balls
model, not only for separately viewed balls.
\end{proof}

In particular, the online outputs at committed vertices are consistent and
give each of them an outgoing edge. To see this formally, extend the query
prefix by all remaining vertices of the current graph. Correctness of $A$
on that full online execution implies correctness and agreement of all the
irrevocable outputs already produced. No future input is needed to run the
simulation.

\begin{lemma}[New commitments are nearby]\label{lem:nearby}
After insertion of $ab$, every newly committed vertex is within distance
$\tau$ of $a$ or $b$ in the updated graph.
\end{lemma}
\begin{proof}
A newly committed vertex $x$ had an old path of length at most $\tau$ to an
old free vertex $f$. If $f$ were still free, the same path would contradict
the new condition $\delta(x)>\tau$. Thus $f$ has ceased to be free, which
can happen only to $a$ or $b$. The old path is still present.
\end{proof}

\subsection{Repair without touching committed edges}

\begin{lemma}[Full-path repair]\label{lem:path}
Let the orientation be legal at every vertex of $C$. Every defect $x\in U$
can be removed by changing directions only on a path of length at most
$\tau$ contained in $U$. This operation creates no new defect and changes
no edge incident to $C$.
\end{lemma}
\begin{proof}
Choose a shortest path
\[
 x=x_0,x_1,\ldots,x_k=f\in F,\qquad k\le\tau.
\]
Every vertex $x_i$ on it has distance at most $k-i\le\tau$ to $F$, and thus
belongs to $U$. Orient the entire path as
\[
 x_0\longrightarrow x_1\longrightarrow\cdots\longrightarrow x_k.
\]
The initial vertex and every internal vertex now have an outgoing path
edge. The last vertex has degree at most two and is exempt. Vertices off
the path have no incident direction changed. Hence no new defect appears,
and the original one disappears. Both endpoints of every changed edge are
in $U$.
\end{proof}

This argument does not require the path to be directed before repair. It
does not require spare outgoing edges at intermediate vertices either: the
new path itself supplies their outgoing edges.

\subsection{The update rule and its radius}

Maintain the following invariant: the orientation is legal, the vertices
already queried from $A$ are exactly $C$, and every edge incident to $C$
has the direction specified by $A$. Initially $C$ is empty and there are no
edges, so the invariant holds.

On insertion of $ab$, perform these steps:
\begin{enumerate}
\item Give $ab$ a provisional direction and recompute $F,C,U$ for the updated
graph. Let $W$ be the entire set of vertices newly entering $C$.
\item Query the single maintained instance of $A$ at all vertices of $W$,
in identifier order. Set their incident directions to the returned values.
Previously committed reports and the new reports agree by
Lemma~\ref{lem:permanence}.
\item While a defect remains in $U$, apply Lemma~\ref{lem:path}, using fixed
tie-breaking rules to choose a defect and a shortest path.
\end{enumerate}
All of $W$ is computed before any repair. If several nearby vertices lose
their last short route to a free vertex, \emph{all} of them are committed
in step~2. Commitment changes orientations, not the degree-based set $F$.

\begin{proof}[Proof of Theorem~\ref{thm:transfer}]
Step~2 makes every vertex of the new $C$ legal and preserves the old
committed directions. Step~3 preserves them and removes at least one defect
without creating another. It terminates and restores the invariant.

It remains to bound all changes during the update. An endpoint of an edge
changed in step~2 is either newly committed or adjacent to a newly committed
vertex. By Lemma~\ref{lem:nearby}, it is within distance $\tau+1$ of $a$ or
$b$. Before step~2, only $a$ or $b$ can have become defective because of
the insertion. After step~2, every defect is therefore within distance
$\tau+1$ of an insertion endpoint.

Every later repair starts at one of these same defects: no repair creates a
new one. Its path has length at most $\tau$. Thus every endpoint of every
changed edge is within distance
\[
 (\tau+1)+\tau=2\tau+1=2r+3
\]
of $a$ or $b$. The repairs do not accumulate radius even when there are
many defects. Inclusion~\eqref{eq:union-intersection} gives the claimed
radius $2r+4$ in the precise convention~\eqref{eq:affected}.
\end{proof}

The proof includes an update that closes an entire component into a cubic
graph. Then that component has no free vertex and all its uncommitted
vertices join $C$. Lemma~\ref{lem:nearby} still places every one of them
near the insertion. No separate closing-component procedure is needed.

\subsection{Vertex insertions}\label{sec:vertices}

To insert a vertex $z$ with its incident edges, first activate its previously
isolated virtual vertex, and then insert its new edges one at a time in
identifier order. Do not expose the intermediate outputs. For a substep
inserting $zv$, the allowed region~\eqref{eq:affected} is contained in
$B(z,R)$ in the substep graph and hence also in the final graph of the
update. The union of all substep changes is therefore within $B(z,R)$.
The same refinement applies to Lemma~\ref{lem:degree-reduction}. Thus vertex
insertions, even with many incident edges in the original graph, do not
increase the radius. An isolated vertex arrival changes no old direction.

\section{Randomization: one static success event for all updates}
\label{sec:random-transfer}

\begin{theorem}[Probability-preserving transfer]\label{thm:random-transfer}
Suppose a randomized \On\ algorithm on subcubic graphs has hard radius
$r(n)\ge1$ and, for every fixed input graph and fixed query order, global
failure probability at most $\varepsilon(n)$. Then there is a randomized
incremental \Dyn\ algorithm with hard radius $2r(n)+4$ and error
$\varepsilon(n)$ against an oblivious insertion adversary.
\end{theorem}
\begin{proof}
Use the deterministic construction with the randomized algorithm $A$ and
one persistent random tape. The sets $F,C,U$, the order within each newly
committed batch, and the refinement of vertex updates all depend only on
the input history. They do not depend on orientations or on random bits.

First specify behavior even on tapes on which $A$ fails. In step~2 collect
all proposed outputs before changing any old direction. Check that they are
well-formed, agree at the two endpoints whenever both are committed, agree
with all old committed directions, and give every newly committed vertex
an outgoing edge. If a check fails, reject the batch, set a permanent failure
flag, and make no old-edge change. From then on, keep all old directions
and orient new edges by a fixed rule. These exceptional executions need not
remain legal, but they always obey the radius bound. If validation succeeds,
the same repair and locality proof as before applies. This makes the
algorithm a total Monte Carlo algorithm with a hard radius on every tape.

Now fix an entire insertion sequence. Let $H_*$ be its final graph on the
$n$ virtual vertices. Let $\sigma$ list, in their scheduled order, every
vertex that ever enters $C$, followed by all other vertices in identifier
order. Both $H_*$ and this \emph{planned} schedule are independent of the
random tape. They are used for the analysis only.

Consider one full execution of $A$ on $(H_*,\sigma)$, using the dynamic
simulation's tape. By Lemma~\ref{lem:permanence}, each input view supplied
by the dynamic simulation is exactly the corresponding view in that fixed
execution. Induction over calls to $A$ shows that the transcripts and
outputs coincide up to any rejected batch. The hypothetical schedule is
defined whether or not the actual simulation aborts.

Let $E$ be the event that this one complete static execution is legal.
Then $\Prb(E)\ge1-\varepsilon(n)$. On $E$, every proposed batch is legal
and compatible with every earlier batch, so no rejection occurs. At every
update, the committed vertices are legal and Lemma~\ref{lem:path} repairs
the uncommitted region. All dynamic outputs are therefore legal on $E$.
This proves the whole-execution probability guarantee with no union bound
over updates.
\end{proof}

\begin{remark}[Scope of the adversary assumption]
If future insertions depend on observed randomized outputs, $H_*$ and
$\sigma$ can depend on the tape, so the fixed-input success theorem for
$A$ does not justify the last argument. Even fixing the final graph while
allowing a coin-dependent insertion order is a different adversary model.
The theorem proved here is precisely the oblivious-sequence statement of
Definition~\ref{def:random}.
\end{remark}

\section{Static tools and the deterministic online algorithm}
\label{sec:static-tools}

We now prove the online bounds used by the transfer. Intermediate graphs
may be loopless multigraphs: parallel edges are distinct edges, count
separately toward degree, and can form a cycle of length two. This causes
no problem for orientation. Identifiers need only be distinct; no bound on
their numerical range is used.

\subsection{Composing phases without looking farther than necessary}

It is convenient to describe a construction as several online phases. Each
phase produces an immutable label at each queried vertex, and a later phase
uses earlier phases' labels as additional input. Global private memory is
permitted in every phase. The following elementary simulation includes the
independence property needed in the randomized proof.

\begin{lemma}[Fixed-cone composition]\label{lem:composition}
On a fixed base graph, suppose there are $p$ online phases with predetermined
radii $\rho_1,\ldots,\rho_p$. A phase can inspect the earlier phases' labels
in its revealed neighborhoods. The final labels can be computed by one
online algorithm of radius $\sum_i\rho_i$. Moreover, for a fixed graph and
fixed external query order, each phase's query order can be chosen
independently of all random bits, even when some vertices are irrelevant
to that phase's task.
\end{lemma}
\begin{proof}
Set $s_i=\sum_{j>i}\rho_j$, so $s_p=0$. On an external query $v$, process
phases $i=1,\ldots,p$ in order. In phase $i$, query every as yet unqueried
vertex of $B(v,s_i)$ in identifier order, storing its returned label.
Irrelevant vertices receive dummy labels but are still included in this
schedule.

If phase $i$ is queried at $u\in B(v,s_i)$, every needed earlier-phase
label lies in $B(v,s_i+\rho_i)$. For each $j<i$,
$s_j\ge s_i+\rho_i$, so phase $j$ has already been processed at all those
vertices. The input required to make this call is within distance
$s_i+\rho_i\le\sum_j\rho_j$ of $v$. The same statement holds for the
unions of previously revealed balls, whose data have been cached.

To verify the simulation, list the calls made to each phase in their actual
order. Imagine first executing phase~1 on its entire list, then phase~2 on
its entire list, and so on. A call sees exactly the same earlier-phase
labels as in the simulation, since all labels it inspects had already been
computed and never change. Its same-phase history and random tape are also
identical. Induction over phases and calls proves equality of the outputs.
After all external vertices are queried, every phase's list includes all
vertices. Finally, the lists were constructed from balls of the fixed base
graph and identifier orders only; they do not depend on label values or
coins.
\end{proof}

This is the online, global-memory version of sequential composition
\cite{GKM,simple}. In particular, the lemma never uses balls in a
coin-dependent residual graph to decide which earlier preprocessing calls
to make.

\subsection{A greedy orientation for high-degree vertices}

\begin{lemma}[High-degree greedy algorithm]\label{lem:high-degree}
Let $Q$ be a loopless multigraph each of whose connected components has at
most $M\ge2$ vertices, and put $h=\lfloor\log_2 M\rfloor+1$. There is a
greedy edge-by-edge orientation, correct for every edge order, which gives
every vertex of degree at least $h$ an outgoing edge. To process an edge,
it needs only its endpoints' degrees, counts of processed incident edges,
and whether they already have outgoing edges.
\end{lemma}
\begin{proof}
Call a vertex \emph{satisfied} if its degree is less than $h$ or an edge has
already been oriented out of it. Process edges once, in any order.
\begin{enumerate}
\item If at least one endpoint is satisfied, orient the edge toward a
satisfied endpoint, with deterministic tie-breaking.
\item Otherwise orient it toward the endpoint with fewer already processed
incident edges, again breaking ties deterministically.
\end{enumerate}
Satisfaction is permanent. Consider the auxiliary graph consisting only of
edges processed by rule~2. Maintain two invariants: each of its components
has at most one unsatisfied vertex, and a component containing an unsatisfied
vertex with $b$ incoming edges has at least $2^b$ vertices.

Initially the assertions hold. Rule~1 cannot give another incoming edge to
an unsatisfied vertex: if one endpoint is unsatisfied, the edge points
\emph{out} of it. Thus this rule can only remove unsatisfied vertices and
does not damage the invariants.

In rule~2, the two unsatisfied endpoints belong to distinct auxiliary
components by the first invariant. Suppose their incoming counts are
$b_u\ge b_v$, and the edge is oriented $u\to v$. Since neither endpoint
has an outgoing edge, these counts equal their processed-edge counts.
The two components merge, $u$ becomes satisfied, and $v$ is the only
possibly unsatisfied vertex of the merged component. Its new incoming
count is $b_v+1$, and the component has at least
\[
 2^{b_u}+2^{b_v}\ge 2^{b_v+1}
\]
vertices. This establishes both invariants.

If a degree-at-least-$h$ vertex were still unsatisfied at the end, all of
its edges would enter it, so its auxiliary component would contain at least
$2^h>M$ vertices. It lies within one connected component of $Q$, a
contradiction.
\end{proof}

The same guarantee applies as soon as all edges incident to a particular
vertex have been processed. Complete the remaining edge order
hypothetically: those incident directions cannot change, and the full
execution is correct by Lemma~\ref{lem:high-degree}.

\subsection{Clusters that either have many boundary edges or can absorb a sink}

Fix an integer $t\ge1$ and set $L=2t+1$. Greedily construct a maximal set
$I$ of centers with pairwise distance greater than $L$: on processing a
vertex, select it if no previously selected center is within distance $L$.
This is an online phase of radius $L$. After all center decisions are made,
assign a vertex $x$ to the center minimizing
\[
 (\dist(x,a),\operatorname{ID}(a))\quad\text{over }a\in I
\]
in lexicographic order. Write $C_a$ for the resulting cluster.

\begin{lemma}\label{lem:clusters}
Every cluster $C_a$ is connected, contains $B(a,t)$, and is contained in
$B(a,L)$. If a cluster contains no cycle and no vertex of original degree
at most two, then its number of boundary edges is at least $2^t$.
\end{lemma}
\begin{proof}
Maximality of $I$ implies that every vertex is within distance $L$ of a
center, proving the radius bound. If $\dist(x,a)\le t$, every other center
$b$ has
$\dist(x,b)\ge\dist(a,b)-\dist(a,x)>2t+1-t>t$,
so $x$ belongs to $C_a$.

For connectedness, let $x$ choose $a$, and let $y$ be the next vertex on a
shortest $x$--$a$ path. If some center $b$ had a better distance--identifier
pair at $y$, the triangle inequality would give a better pair for $b$ at
$x$ too. Thus $y$ also chooses $a$, and repeating the argument gives a
path inside $C_a$ from $x$ to $a$.

Now suppose $C_a$ is acyclic and all its vertices have original degree at
least three. Its induced subgraph is a tree. Its contained radius-$t$ ball
has at least $3\cdot2^{t-1}$ vertices at distance $t$ from $a$: the root
has at least three children, and every vertex at a smaller positive depth
has at least two children. A missing or identified child would contradict
either full degree in that ball or acyclicity. In particular
$|C_a|\ge2^t$. If $b(C_a)$ counts its boundary edges with multiplicity,
the degree sum identity gives
\[
 b(C_a)=\sum_{x\in C_a}\deg_G(x)-2(|C_a|-1)
 \ge |C_a|+2\ge2^t.
\]
\end{proof}

Call a cluster \emph{absorbing} if it contains a degree-at-most-two vertex
or a cycle. Its internal edges can be oriented so that all its constrained
vertices have outgoing edges regardless of boundary directions. For a free
vertex, orient a spanning tree toward it. For a cycle, orient that cycle
cyclically and orient a forest spanning the remaining vertices toward the
cycle. Orient unused edges arbitrarily. All choices use identifiers and
edge identities, so every vertex can reconstruct the same orientation.

More generally, if a cluster has an outgoing boundary edge at $w$, orient a
spanning tree toward $w$. Then $w$ is satisfied by that boundary edge and
all other cluster vertices are satisfied by the tree. These observations
apply to a two-edge cycle in a multigraph as well.

\subsection{The deterministic construction and its online implementation}

\begin{theorem}[Componentwise deterministic static bound]\label{thm:static-det}
Let $G$ be a loopless multigraph each of whose connected components has at
most $M\ge2$ vertices. Sinkless orientation can be computed by a deterministic
online algorithm of radius $O(\log\log M)$. It needs only the bound $M$,
not component sizes or the number of components. If a component exceeds
$M$, the same algorithm still returns a consistent orientation within the
same radius, although that component may contain defects.
\end{theorem}
\begin{proof}
Put
\[
 h=\lfloor\log_2 M\rfloor+1,\qquad
 t=\lceil\log_2 h\rceil,\qquad L=2t+1.
\]
Construct the clusters above. Contract each cluster to a vertex, keeping
every intercluster edge as a distinct edge of a loopless quotient multigraph
$Q$. Each component of $Q$ has at most $M$ vertices. Apply
Lemma~\ref{lem:high-degree} to its edges.

A quotient vertex of degree at least $h$ has an outgoing edge. A cluster
with fewer than $h$ boundary edges must be absorbing by
Lemma~\ref{lem:clusters}, since $2^t\ge h$. Thus every cluster either has
an outgoing boundary edge or is absorbing. Orient its interior as described
above. The resulting orientation is sinkless.

Here is a complete local online implementation, avoiding any assumption
that the full quotient graph is already known. Use three phases on $G$.
\begin{enumerate}
\item Compute the greedy center set with radius $\rho_1=L$.
\item At each vertex $v$, compute and store the full description of its
cluster, including its induced edges, original vertex degrees, boundary
edges, and the owners of their other endpoints. Radius $\rho_2=3L+3$
in the graph labeled with the completed center decisions suffices.
Indeed, $v$'s center is within $L$, its entire cluster is within $2L$
of $v$, and owners of boundary neighbors can be determined from centers
within a further $L+1$ hops. The spare layers provide every needed
incidence and degree.
\item On the first query at a member $v$ of a cluster $C$, process every
still unoriented boundary edge of $C$ by the high-degree greedy rule, in
fixed order. Store the directions globally. For each adjacent cluster,
its complete phase-2 description is available at the other endpoint of
the boundary edge, which lies within $2L+1$ of $v$. Degrees, counts of
previously processed edges, and outgoing-edge status of quotient vertices
can therefore be obtained from these descriptions and the persistent
greedy history. Then orient and store the entire interior of $C$.
Return the directions at $v$. On a later query in $C$, return the stored
directions. This phase has radius $\rho_3=2L+2$.
\end{enumerate}
Processing a cluster finishes all of its boundary edges, so neither their
directions nor its chosen interior ever changes afterward. Each quotient
edge is processed exactly once. The resulting edge order can be arbitrary;
Lemma~\ref{lem:high-degree} applies to all orders. Returning all incident
directions therefore gives consistent immutable endpoint outputs.

By Lemma~\ref{lem:composition}, the three phases require radius at most
$6L+5=O(\log\log M)$ in the original graph. The argument uses the bound
$M$ only in the high-degree lemma, component by component. If the bound is
false, make the algorithm total as follows: whenever a cluster has neither
an outgoing boundary edge nor an absorbing structure, orient its spanning
tree toward its least-identifier vertex and orient the rest arbitrarily.
This may leave one defect, but produces a consistent orientation by the
same local computations. Thus the radius bound is unconditional.
\end{proof}

In particular, taking $M=n$ gives the deterministic online bound required
by Theorem~\ref{thm:transfer}. The componentwise formulation will also let
us solve many small randomized residual components without exploring an
entire component first.

\section{The randomized online algorithm, with a full shattering proof}
\label{sec:static-random}

\begin{theorem}[Randomized static bound]\label{thm:static-random}
For every fixed $c>0$, sinkless orientation on a subcubic graph with a known
vertex bound $n\ge2$ has a randomized online algorithm of radius
$O_c(\log\log\log n)$ and global error at most $n^{-c}$ for every fixed
graph and fixed query order. The radius holds on every random tape.
\end{theorem}

The proof has three combinatorial steps: replace constant-radius clusters
by vertices with many available incidences, randomly satisfy most of those
vertices, and use Theorem~\ref{thm:static-det} on the remaining small
components. We then explain the complete online implementation. This is the
marking approach of~\cite{GS,GHK}, with the reductions and probability
estimate given here explicitly.

\subsection{Constant-radius clusters and half-edges}

Fix the absolute constants
\[
 a=8,\qquad D=512.
\]
Greedily choose a maximal set of centers at pairwise distance greater than
$2a=16$, and assign every vertex to its nearest center, breaking ties by
center identifier. The proof of Lemma~\ref{lem:clusters} applies with this
separation too: clusters are connected, have radius at most $16$ and strong
diameter at most $32$, and contain the full radius-$8$ ball of their center.
Indeed, a different center is at distance at least $17-8=9$ from a vertex
within distance $8$ of the chosen center.

As before, a cluster is absorbing if it contains a cycle or a vertex of
degree at most two. Fix an internal orientation satisfying all its
degree-three vertices, independent of its boundary. Call every other
cluster \emph{active}. An active cluster $C$ is a tree whose vertices all
have original degree three. Its radius-$8$ ball is a full cubic tree, and
therefore
\begin{equation}\label{eq:boost}
 |C|\ge1+3\sum_{i=0}^{7}2^i=766,\qquad
 b(C)=3|C|-2(|C|-1)=|C|+2\ge768>D.
\end{equation}

At each active cluster select exactly $D$ boundary edges by a fixed order
on their original endpoint identifiers. Form a quotient instance $Q$ with
one vertex per active cluster. An original boundary edge selected at both
endpoints becomes an ordinary quotient edge. If selected at exactly one
endpoint, it becomes a \emph{half-edge} at that endpoint. If selected at
neither endpoint, discard it from $Q$ and fix an arbitrary direction for it.
An absorbing endpoint never selects an edge.

A half-edge has one represented endpoint, and can be directed inward or
outward relative to it. An outward half-edge satisfies that endpoint. Each
item of $Q$, ordinary edge or half-edge, corresponds to a distinct original
boundary edge. Thus any orientation of all the items maps consistently
back to original edges. Every quotient vertex has exactly $D$ incident
items, counted with multiplicity. The ordinary underlying graph has maximum
degree at most $D$; parallel edges are permitted, and there are no loops.

If every active quotient vertex has an outgoing item, every active original
cluster has an outgoing boundary edge. Orient its internal tree toward the
tail of one such edge. Absorbing clusters use their fixed internal solutions.
This gives a sinkless orientation of the original graph. An unrepresented
endpoint of a half-edge causes no conflict: it must obtain its own witness
from its own selected items, or it belongs to an absorbing cluster. It places
no requirement on that particular edge.

All of this preprocessing has constant radius. A step in the quotient
corresponds to travel within a cluster and across one original boundary
edge, costing at most $32+1=33$ original hops, with an additive constant
to inspect endpoint clusters. We will ensure below that the entire
preprocessing is independent of the later random bits.

\subsection{Marking and a residual orientation problem}

Independently for every item of $Q$, mark it with probability $1/4$. Give
each marked item a uniformly random direction. There is one common random
choice per ordinary edge, not separate choices at its endpoints; for a
half-edge the choices are inward and outward.

Let $X_v$ count the marked items incident to a quotient vertex $v$, and let
$Z_v$ be the event that $v$ has no outgoing marked item. Classify bad vertices
in the following priority order:
\begin{enumerate}
\item \emph{Type I:} $X_v>D/2$.
\item \emph{Type II:} not type I, but adjacent by an ordinary edge to a
type I vertex.
\item \emph{Type III:} neither type I nor type II, and $Z_v$ occurs.
\end{enumerate}
All remaining vertices are good. Let $B$ be the set of bad vertices. Unmark
every item incident to a type I vertex, and freeze the directions of all
marks that remain. For a half-edge, only its represented endpoint counts
as incident; the other original cluster has no role in this rule.

\begin{lemma}\label{lem:marking}
Every good vertex has an outgoing marked item that remains frozen.
Every bad vertex has at least $D/2$ unmarked incident items.
\end{lemma}
\begin{proof}
A good vertex has an outgoing marked item because it is not type III, and
no such item is unmarked: neither the vertex nor any ordinary neighbor is
type I. A type I vertex has all its $D$ items unmarked. Any other bad vertex
initially had at most $D/2$ marks, and unmarking only increases the number
of unmarked items.
\end{proof}

Keep unmarked edges between bad vertices as residual ordinary edges, and
replace an unmarked edge between a bad and a good vertex by a residual
half-edge at the bad vertex. An unmarked original half-edge at a bad vertex
remains a half-edge. Orient all other unmarked items arbitrarily. The
residual instance has vertex set $B$, with at least $D/2=256$ incident
items at each vertex. A sinkless orientation of this instance finishes the
orientation of $Q$: every bad vertex gets an outgoing item, while no good
vertex loses its frozen outgoing witness.

\subsection{Why every residual component is small with high probability}

We prove the probability bound instead of invoking a shattering theorem.
All distances and neighborhoods in this subsection use the ordinary
underlying graph of $Q$ and ignore half-edges.

Define a \emph{primitive rare event}
\[
 P_v=\{X_v>D/2\}\cup Z_v.
\]
It depends only on the independent random choices for items incident to
$v$. Consequently these events are mutually independent at any set of
pairwise nonadjacent vertices. This remains true for parallel edges and
half-edges. We use primitive events rather than literal type III events,
since the latter also inspect neighbors' type I status.

An exponential-moment estimate and independence of the $D$ incident items
give
\begin{align}
 \Prb[X_v>D/2]
 &\le 2^{-D/2}\EE[2^{X_v}]
  =2^{-D/2}(5/4)^D\le e^{-D/12},\label{eq:mark-tail}\\
 \Prb[Z_v]&=(7/8)^D\le e^{-D/8}.\label{eq:zero-tail}
\end{align}
For the last inequality in~\eqref{eq:mark-tail}, it suffices to use
$\ln(5/4)\le1/4$ and $\ln2\ge2/3$. The latter follows, for example,
by integrating the convex function $1/x$ on $[1,2]$ and using its midpoint
value. Thus
\begin{equation}\label{eq:rare-prob}
 \Prb[P_v]\le p:=2e^{-D/12}.
\end{equation}

Let $P$ be the set of vertices where $P_v$ occurs, and let $I$ be the type I
set. The definitions imply
\begin{equation}\label{eq:bad-cover}
 B=P\cup N_Q(I).
\end{equation}
In particular, each bad vertex is primitive or adjacent to a primitive
vertex, and such a witnessing neighbor lies in the same bad component.

\begin{lemma}[Separated witnesses]\label{lem:witness}
If a component of $Q[B]$ has at least $2D^2k$ vertices, there is a set of
$k$ primitive vertices that are pairwise nonadjacent in $Q$ and connected
by pairs whose $Q$-distance is between two and four.
\end{lemma}
\begin{proof}
Choose any primitive vertex in the component, which exists by
\eqref{eq:bad-cover}, and start a set $S$ with it. A radius-two ball in a
graph of maximum degree at most $D$ has at most
$1+D+D(D-1)=D^2+1\le2D^2$ vertices. While $|S|<k$, the radius-two
neighborhood of $S$ therefore does not cover the component.

Take a path inside the component from $S$ to an uncovered vertex. Distances
to $S$ change by at most one along an edge, so this path contains a vertex
$w$ at $Q$-distance exactly three from $S$. If $w$ is primitive, add it to
$S$. Otherwise, by~\eqref{eq:bad-cover}, it has a type I neighbor $z$ in
the same component. Add $z$ instead. In the latter case
$2\le\dist_Q(S,z)\le4$; in the former case the distance is three.
The new vertex is nonadjacent to every previous choice and within distance
four of at least one. Induction gives the required independent, connected
witness set.
\end{proof}

Let $J$ connect pairs of vertices at $Q$-distance two, three, or four. Its
maximum degree is at most
\[
 D^2+D^3+D^4\le2D^4=:\Lambda.
\]
There are at most $n(4\Lambda)^{k-1}$ rooted ordered trees with $k$ distinct
vertices embedded in $J$. Indeed, at most $4^{k-1}$ tree shapes are possible:
a depth-first traversal is encoded by $k-1$ opening and $k-1$ closing steps.
There are at most $n$ choices for the root and at most $\Lambda$ choices
for the image of each child. Restricting to injective embeddings whose
images are independent in $Q$ only decreases this count.

For any such independent image, all $k$ primitive events occur with
probability at most $p^k$. Every witness from Lemma~\ref{lem:witness} has
a spanning tree counted above. A union bound therefore gives
\begin{equation}\label{eq:witness-prob}
 \Prb[\text{a primitive witness of size }k\text{ exists}]
 \le n(4\Lambda)^{k-1}p^k\le n(4\Lambda p)^k.
\end{equation}
The fixed choice $D=512=2^9$ satisfies
\[
 4\Lambda p\le16D^4e^{-D/12}
 =\exp(40\ln2-128/3)<e^{-2}.
\]
For the chosen confidence exponent $c>0$, set
\begin{equation}\label{eq:residual-bound}
 k=\lceil(c+1)\ln n\rceil,\qquad m_0=2D^2k.
\end{equation}
Equations~\eqref{eq:witness-prob} and~\eqref{eq:residual-bound} show that
the probability of a bad component with at least $m_0$ vertices is at most
\begin{equation}\label{eq:shatter-success}
 ne^{-2k}\le n^{-2c-1}\le n^{-c}.
\end{equation}
The residual ordinary graph only deletes edges from $Q[B]$. Thus with
probability at least $1-n^{-c}$, \emph{all} its components have fewer than
$m_0=O_c(\log n)$ vertices. Call this the shattering event.

\subsection{Cleaning up residual half-edges}

Replace each residual half-edge by an ordinary edge to its own fresh
degree-one leaf. The resulting graph $R$ is a loopless multigraph. Old
residual vertices have degree at least $256$, while the leaves are exempt
from the sinkless requirement. A direction on a leaf edge translates
directly to a direction on its half-edge.

There are at most $D$ such leaves per residual vertex. On the shattering
event, every component of $R$ consequently has at most
\begin{equation}\label{eq:cleanup-M}
 M=(D+1)m_0=O_c(\log n)
\end{equation}
vertices. Give all virtual vertices distinct tuple identifiers, derived
from their cluster identifiers and, for leaves, their original edge
identities, with separate type tags for quotient vertices and leaves.
Apply Theorem~\ref{thm:static-det} to $R$ with the numerical
parameter $M$. It gives each old residual vertex an outgoing edge, using
the explicit radius from that theorem's proof,
\begin{equation}\label{eq:cleanup-radius}
 q=12\left\lceil\log_2\bigl(\lfloor\log_2 M\rfloor+1\bigr)\right\rceil+11
   =O_c(\log\log\log n).
\end{equation}
No routine first explores or counts a residual component. The algorithm
uses this same bound $M$ in every component, regardless of the total number
of virtual vertices or the range of their identifiers. On tapes where a
component exceeds $M$, its off-promise behavior still uses radius $q$ and
returns consistent directions, by Theorem~\ref{thm:static-det}.

\subsection{A fully local online simulation}

It remains to justify both online locality and the independence assumed in
the probability calculation. In particular, an online algorithm cannot
silently assume that all clusters were computed globally in advance.

Use Lemma~\ref{lem:composition} on the original graph $G$ with the following
six phases. Each phase returns an immutable label at \emph{every} original
vertex, using a dummy label when its particular task is irrelevant there.
The numerical radii below are deliberately loose constants.
\begin{enumerate}
\item \emph{Centers}, radius $16$: compute the greedy set of centers at
distance greater than $16$.
\item \emph{Owners}, radius $16$: read the completed center labels and
choose the nearest center, breaking ties by identifier.
\item \emph{Cluster descriptions}, radius $100$: store at each member its
cluster's entire induced graph, degrees, boundary edges, absorbing status,
and selected $D$ incidences if active. A member is within distance $32$
of every vertex in its cluster, and its boundary neighbors are within
distance $33$, so the stated radius more than suffices. Store a canonical
internal solution for each absorbing cluster.
\item \emph{Random marking}, radius $100$: construct the selected quotient
incidences, generate their marks and directions, and record types and
surviving marks. To inspect an ordinary neighbor's type I event it suffices
to inspect that neighboring cluster's selected incidences. Two cluster
steps and their complete endpoint data fit within this radius, using the
phase-3 descriptions. Generate each item's random choice once, indexed by
its original edge identity, and cache it. Thus two endpoints always use
the same choice.
\item \emph{Residual cleanup}, radius $100(q+3)$: simulate the deterministic
algorithm on $R$. When this phase first processes any member of a bad
active cluster, query the corresponding residual vertex and all its
private leaf vertices, in a fixed order, and cache their outputs.
Later members reuse these outputs. A virtual step projects to at most
$33$ original hops, and a leaf stays at its owner's cluster. The
radius-$q$ virtual views for this block, including the needed incidence
data and bad/good labels, are therefore within the stated original
radius. The phase returns its cluster's complete selected-item directions,
combining surviving marks with all residual outputs, including directions
on private-leaf edges. At a good active cluster, the label instead records
its marked witnesses and the fixed directions already determined there;
directions controlled by a bad neighbor are read from that neighbor in
the final phase.
\item \emph{Final orientations}, radius $100$: reconstruct the direction
of every original boundary edge from the earlier labels at its selected
endpoint or endpoints. Orient each active cluster's internal tree toward
the tail of a selected outgoing item; absorbing clusters use their stored
solutions. Return all incident directions of the queried vertex. If an
active cluster has no outgoing item on an unsuccessful tape, orient its
tree toward its least-identifier vertex instead. All members use the same
complete data and tie-breaking rules, so their reports agree.
\end{enumerate}

For clarity, a phase-5 label at a good cluster need not predict the output
of a neighboring bad cluster. Phase~6 has the completed earlier-phase
labels throughout its fixed-radius view, so it obtains that direction at
the bad endpoint. Likewise, private leaf queries use the same global
instance of the componentwise deterministic algorithm and give consistent
directions at both ends of every virtual edge. All virtual vertices are
eventually queried when their original clusters are processed.

Each phase is a valid online procedure with the prescribed radius on its
earlier-phase-labeled input. By Lemma~\ref{lem:composition}, the total
original radius is at most $100q+632$, hence is
$O_c(\log\log\log n)$. This includes all preprocessing and all lifting
overheads, even on unsuccessful tapes.

Fix now the original input graph and its complete external query order.
The cone simulation fixes the query order of every preprocessing phase
independently of all random bits. Consequently the completed centers,
clusters, selected incidences, and quotient $Q$ are fixed independently of
the marking coins. This is exactly the hypothesis needed for
\eqref{eq:rare-prob}--\eqref{eq:shatter-success}. The virtual residual query
order may depend on which clusters are bad, but that causes no difficulty:
the deterministic cleanup works for every order whenever the component
bound holds.

\begin{proof}[Proof of Theorem~\ref{thm:static-random}]
The serialized execution of the six phases is correct on the shattering
event: Lemma~\ref{lem:marking} protects good quotient vertices, the
componentwise cleanup satisfies all bad ones, and the cluster lifting
satisfies every original degree-three vertex. The online simulation
produces exactly those outputs by Lemma~\ref{lem:composition}. The
shattering event has probability at least $1-n^{-c}$, and the radius bound
holds unconditionally, as just established.
\end{proof}

\section{Completing the upper bounds}\label{sec:conclude}

\begin{proof}[Proof of Theorem~\ref{thm:main}]
For the deterministic result, apply Theorem~\ref{thm:static-det} with
$M=n$ to obtain a subcubic online algorithm of radius $O(\log\log n)$.
Theorem~\ref{thm:transfer} gives the same asymptotic dynamic radius.
Lemma~\ref{lem:degree-reduction} extends it to general graphs, and
Section~\ref{sec:vertices} handles vertex insertions with incident edges.

For the randomized result, Theorem~\ref{thm:static-random} supplies an
online radius $O_c(\log\log\log n)$ and error $n^{-c}$ on every fixed
subcubic graph and fixed query order. Apply
Theorem~\ref{thm:random-transfer}, then the same degree reduction and
vertex-update refinement. For a fixed original insertion sequence all
these reductions are deterministic. The one static success event in the
transfer guarantees legal orientations at every dynamic time, with no
additional probability loss.
\end{proof}

The transfer relies on both bounded-degree saturation and the freedom to
leave low-degree vertices as sinks. It is therefore a statement about this
particular orientation problem, rather than a general equivalence of the
models. The proof also relies on insertions only: deletions could reopen a
previously permanent neighborhood.

Neither upper bound is asserted to be optimal. For context, the lower
bounds of Brandt and G\"ottlicher~\cite[Theorem 1.2 and Corollary 1.3]{lower},
together with the dynamic-to-online simulation of~\cite{models}, give a
$2^{\Omega(\log^* n)}$ barrier for both deterministic dynamic locality and
the oblivious randomized version considered here. Those lower bounds are
background results, not ingredients in the self-contained upper-bound
proofs. Determining the optimal locality within the remaining gaps is open.

\begin{thebibliography}{99}
\bibitem{models}
Amirreza Akbari, Navid Eslami, Henrik Lievonen, Darya Melnyk,
Joona S\"arkij\"arvi, and Jukka Suomela.
\newblock Locality in Online, Dynamic, Sequential, and Distributed Graph Algorithms.
\newblock In \emph{ICALP 2023}, LIPIcs 261, 10:1--10:20, 2023.
\newblock \href{https://doi.org/10.4230/LIPIcs.ICALP.2023.10}{doi:10.4230/LIPIcs.ICALP.2023.10}.
\newblock Full version: \url{https://arxiv.org/abs/2109.06593v4}.

\bibitem{simple}
Alkida Balliu, Janne H. Korhonen, Fabian Kuhn, Henrik Lievonen,
Dennis Olivetti, Shreyas Pai, Ami Paz, Joel Rybicki, Stefan Schmid,
Jan Studen\'y, Jukka Suomela, and Jara Uitto.
\newblock Sinkless Orientation Made Simple.
\newblock In \emph{SOSA 2023}, pages 175--191, 2023.
\newblock \href{https://doi.org/10.1137/1.9781611977585.ch17}{doi:10.1137/1.9781611977585.ch17}.
\newblock Full version: \url{https://arxiv.org/abs/2108.02655v2}.

\bibitem{GHK}
Mohsen Ghaffari, David G. Harris, and Fabian Kuhn.
\newblock On Derandomizing Local Distributed Algorithms.
\newblock In \emph{FOCS 2018}, pages 662--673, 2018.
\newblock \href{https://doi.org/10.1109/FOCS.2018.00069}{doi:10.1109/FOCS.2018.00069}.
\newblock Full version, 17 September 2019:
\url{https://arxiv.org/abs/1711.02194v4}.
\newblock Theorem numbers in this note refer to that full version.

\bibitem{GS}
Mohsen Ghaffari and Hsin-Hao Su.
\newblock Distributed Degree Splitting, Edge Coloring, and Orientations.
\newblock In \emph{SODA 2017}, pages 2505--2523, 2017.
\newblock \href{https://doi.org/10.1137/1.9781611974782.166}{doi:10.1137/1.9781611974782.166}.
\newblock Full version: \url{https://arxiv.org/abs/1608.03220}.

\bibitem{GKM}
Mohsen Ghaffari, Fabian Kuhn, and Yannic Maus.
\newblock On the Complexity of Local Distributed Graph Problems.
\newblock In \emph{STOC 2017}, pages 784--797, 2017.
\newblock \href{https://doi.org/10.1145/3055399.3055471}{doi:10.1145/3055399.3055471}.
\newblock Full version: \url{https://arxiv.org/abs/1611.02663v2}.

\bibitem{mending}
Alkida Balliu, Juho Hirvonen, Darya Melnyk, Dennis Olivetti,
Joel Rybicki, and Jukka Suomela.
\newblock Local Mending.
\newblock In \emph{SIROCCO 2022}, LNCS 13298, pages 1--20, 2022.
\newblock \href{https://doi.org/10.1007/978-3-031-09993-9_1}{doi:10.1007/978-3-031-09993-9\_1}.
\newblock Full version: \url{https://arxiv.org/abs/2102.08703v4}.

\bibitem{coloring}
Yann Bourreau, Filippo Casagrande, Tim G\"ottlicher, and
Ananth Narayanan Karunattillam.
\newblock Dynamic 3-Coloring of Bipartite Graphs with Optimal Locality.
\newblock Preprint, 14 October 2026.
\newblock \url{https://arxiv.org/abs/2610.19710v1}.

\bibitem{quantum}
Amirreza Akbari, Xavier Coiteux-Roy, Francesco d'Amore, Fran\c{c}ois Le Gall,
Henrik Lievonen, Darya Melnyk, Augusto Modanese, Shreyas Pai,
Marc-Olivier Renou, V\'aclav Rozho\v{n}, and Jukka Suomela.
\newblock Online Locality Meets Distributed Quantum Computing.
\newblock Conference version in \emph{STOC 2025}.
\newblock \href{https://doi.org/10.1145/3717823.3718211}{doi:10.1145/3717823.3718211}.
\newblock We refer to the corrected full version of 13 February 2026:
\url{https://arxiv.org/abs/2403.01903v4}.

\bibitem{lower}
Sebastian Brandt and Tim G\"ottlicher.
\newblock A Post-Quantum Lower Bound for the Distributed Lov\'asz Local Lemma.
\newblock In \emph{SODA 2026}, 2026.
\newblock \href{https://doi.org/10.1137/1.9781611978971.119}{doi:10.1137/1.9781611978971.119}.
\newblock Full version: \url{https://arxiv.org/abs/2510.15698v1}.

\bibitem{Nolin}
Alexandre Nolin.
\newblock Brief Announcement: Sinkless Orientation Made Trivial.
\newblock In \emph{PODC 2026}, pages 333--336, 2026.
\newblock \href{https://doi.org/10.1145/3796701.3815953}{doi:10.1145/3796701.3815953}.
\newblock Author version: \url{https://hal.science/hal-05678283v1}.
\end{thebibliography}

\end{document}
